\section{Dữ liệu, tài khoản, triển khai và chi phí vận hành}\label{sec:data-design}
Dữ liệu cốt lõi là một sổ sự kiện chỉ ghi thêm cho mỗi buổi học. Mỗi sự kiện có kiểu, nội dung, thời điểm, số thứ tự trong buổi và mã trường; các kiểu gồm phản hồi học sinh, lời gia sư, quyết định chính sách, câu hiển thị, câu trả lời đã chấm, gợi ý, câu được khép, lệnh gọi và chuyển mô hình, phản hồi bị từ chối và phần thưởng bandit. Chính sách, điểm hiểu bài và mọi báo cáo là phép chiếu tất định từ sổ này, nên phát lại một buổi không cần gọi lại mô hình: lựa chọn của mô hình đã nằm trong sổ. Cơ sở dữ liệu là MongoDB Atlas gói miễn phí M0, không hỗ trợ giao dịch nhiều tài liệu, nên trạng thái và sự kiện của một buổi nằm trong một tài liệu ghi nguyên tử với số phiên bản để phát hiện hai lần ghi chồng nhau. Tệp gốc và âm thanh nằm trên ổ đĩa máy chủ. Database phát triển, kiểm thử và dữ liệu thật được tách tuyệt đối: kiểm thử chạm cơ sở dữ liệu thật chỉ chạy trên database tạm có tên theo mẫu cho phép và tự dọn, dữ liệu demo chỉ tạo bằng một lệnh riêng và gắn cờ để xóa được, và không script nào tự ghi vào database phát triển hay thật.

Giáo viên đăng nhập bằng mã PIN lưu dạng băm một chiều. Học sinh tự tạo tài khoản bằng tên đăng nhập và mã bí mật bốn đến tám chữ số. Phụ huynh dùng chung tài khoản của con và nhận mã vai trò riêng nên chỉ xem được báo cáo của con. Máy chủ ký mã xác thực bằng khóa riêng của lần triển khai, mọi ràng buộc quyền được thực thi ở máy chủ theo vai trò, và máy khách không có đường nào chọn mô hình hay vượt ràng buộc của buổi học.

Khi phát triển, hệ thống chạy trực tiếp trên máy không dùng Docker, một cổng phục vụ cả web và API. Khi lên máy chủ, người vận hành dùng Docker Compose với một ảnh chứa web và API cùng một MongoDB riêng có mật khẩu và không mở cổng ra ngoài. Dữ liệu database nạp từ tệp xuất bằng lệnh chỉ đọc; vì kho mã công khai và tệp chứa lời các em đã nói, tệp này không bao giờ nằm trong kho mà đi qua kênh riêng. Bộ kiểm thử chạy hoàn toàn ngoại tuyến, không cần mạng hay database thật.

Về chi phí và khả năng mở rộng, hệ thống chạy hoàn toàn trên dịch vụ miễn phí nên chi phí tiền bằng không, nhưng hạn mức là giới hạn thật. Mỗi lượt gọi mô hình vài lần, mỗi lần tốn trung bình khoảng $1\,100$ token và lời nhắc không quá $3\,000$ token, nên hạn mức khoảng $8\,000$ token mỗi phút giới hạn số học sinh học đồng thời ở mức vài em trên một khóa; trần số lệnh gọi theo buổi và theo chương giữ cho một buổi không dùng hết hạn mức của người khác, và câu mẫu đảm bảo buổi học không dừng khi hạn mức cạn. Đồ án chưa đo độ trễ từng lượt và chưa kiểm chứng việc phục vụ nhiều học sinh cùng lúc; việc mở rộng đòi chuyển sang gói trả phí hoặc mô hình chạy cục bộ, và đây là một hạn chế được nêu ở mục Kết luận.

\section{Tổng kết chương}
Chương này trình bày thiết kế trong đó hệ thống học tập được dựng trước và mô hình ngôn ngữ chỉ giữ vai diễn đạt. Sáu bất biến ở mục~\ref{sec:nfr} được giữ bằng các thành phần tất định: vòng sáu bước và thang gợi ý quyết định điều gia sư làm, truy xuất lai theo đồ thị giới hạn điều gia sư được dạy, bộ chấm và phép chiếu mức hiểu quyết định điều học sinh đã làm được, và bộ kiểm chứng chặn lời của mô hình khi vi phạm. Bộ chọn bandit (mục~\ref{sec:bandit-design}) là phần học được duy nhất và vẫn nằm trong thực đơn mà chính sách cho phép. Chương sau đánh giá các thành phần này.
